\documentclass[11pt]{article}

\usepackage[a4paper, margin=1in]{geometry}
\usepackage{graphicx}
\usepackage{booktabs}
\usepackage{amsmath}
\usepackage{hyperref}
\usepackage{float}
\usepackage{enumitem}
\usepackage{parskip}
\usepackage{microtype}

\hypersetup{
  colorlinks=true,
  linkcolor=blue,
  urlcolor=blue,
  citecolor=blue
}

\title{Detecting Patronising and Condescending Language (PCL) \\
Natural Language Processing Coursework 2026}
\author{Darius Aron \\
Leaderboard Name: \texttt{dodo} \\
GitHub Repository: \url{https://github.com/darius024/pcl-classifier}}
\date{}

\begin{document}

\maketitle
\thispagestyle{empty}

\newpage
\tableofcontents

% ============================================================
\clearpage
\section{Exercise 1: Critical Paper Review}
\label{sec:review}

\subsection{Q1: Primary Contributions}

The paper ``Don't Patronize Me! An Annotated Dataset with Patronizing and Condescending Language towards Vulnerable Communities'' by Pérez-Almendros et al.\ [1, 2] makes three primary contributions:

\begin{itemize}
    \item \textbf{Dataset:} The authors introduce the \emph{Don't Patronize Me!} (DPM) corpus, consisting of over 10,000 paragraphs extracted from English-language news articles, each manually annotated for the presence of PCL towards vulnerable communities. Critically, annotations are made at both the document level (binary PCL/No-PCL) and the text-span level, where annotators highlight the specific spans that contain PCL, enabling fine-grained supervision at multiple levels. The paragraphs were collected by keyword search (e.g., ``homeless'', ``refugee'', ``disabled'') across articles from many countries, providing broad coverage.

    \item \textbf{Taxonomy:} The paper proposes a two-level, hierarchical taxonomy of PCL types focused on language directed at vulnerable communities. The first level distinguishes PCL from non-PCL. The second level decomposes PCL into seven fine-grained categories: \emph{Unbalanced Power Relations}, \emph{Shallow Solution}, \emph{Presupposition}, \emph{Authority Voice}, \emph{Metaphor}, \emph{Compassion}, and \emph{The Poorer the Merrier}. This taxonomy brings conceptual clarity to a previously informal notion, making annotation reproducible and enabling multi-label classification (Subtask 2 of SemEval 2022).

    \item \textbf{Community benchmark:} Beyond organising the competition, the authors [2] established a reusable benchmark by releasing the corpus with official train/dev/test splits, evaluation scripts, and baseline results across several transformer architectures (including RoBERTa-base at F1~=~0.48). This gives subsequent researchers a concrete foundation: new approaches can be compared against documented systems on identical data without re-collecting or re-annotating, lowering the barrier for reproducible follow-on work.
\end{itemize}

\subsection{Q2: Technical Strengths}

\begin{itemize}
    \item \textbf{Originality and domain novelty:} PCL directed at vulnerable communities had received almost no prior computational treatment. This work fills a clear gap by operationalising a socially important phenomenon that sits between well-studied domains such as hate speech and sarcasm detection. The taxonomy itself is a standalone scholarly contribution that could inform future annotation guidelines and datasets.

    \item \textbf{Annotation rigour:} The authors employed a multi-annotator setup with trained annotators, reported inter-annotator agreement (IAA) using standard metrics (Fleiss's $\kappa$), and used a majority-voting scheme to produce the final binary labels. The annotation guidelines were developed iteratively through calibration rounds, which increases the reliability of gold labels and reduces the risk of systematic annotator bias.

    \item \textbf{Evaluation design:} The choice of F1 score on the positive (PCL) class as the primary metric is well-motivated given the strong class imbalance (roughly 10\% PCL). Accuracy would be misleading: a model that always predicts ``No PCL'' achieves $\sim$90\% accuracy with no useful signal. F1 on the minority class directly captures both precision and recall for the phenomenon of interest.
\end{itemize}

\subsection{Q3: Weaknesses}

\begin{itemize}
    \item \textbf{Low inter-annotator agreement:} The paper reports low IAA scores, attributed to the inherent subjectivity of PCL. While this is acknowledged, the analysis is superficial. There is no investigation into \emph{which} PCL categories or \emph{which} severity levels drive the most disagreement, nor whether specific annotator pairs systematically diverge. A deeper breakdown would clarify whether the low agreement reflects genuine ambiguity in the phenomenon or inconsistencies in the guidelines that could be addressed. More annotators or a different elicitation process (e.g., pairwise comparison rather than independent labelling) might have yielded higher and more stable agreement.

    \item \textbf{Shallow model analysis:} The baseline models presented (including RoBERTa-base) are described without detailed justification for hyperparameter choices such as learning rate, batch size, or class weighting strategy. There are no ablation studies showing the contribution of individual design decisions. Without such analysis, it is hard to tell whether the reported F1 of 0.48 is a reasonable lower bound or whether better tuning of the same architecture could have yielded significantly higher results.

    \item \textbf{Limited generalisation discussion:} The dataset is drawn exclusively from English-language news articles via keyword-based search, meaning all examples are topically anchored to the chosen vulnerable-community terms. The authors do not discuss whether the taxonomy and models would generalise to other domains (e.g., social media, parliamentary speech, or informal text), other languages, or PCL targeting groups not included in the keyword list. The keyword-seeded collection strategy also risks introducing spurious correlations between keywords and labels that could inflate model performance without reflecting genuine PCL understanding. A related concern is uneven keyword representation: some PCL categories (for example, \textit{The Poorer the Merrier}) appear far more frequently than others, creating a within-class imbalance. This may cause models to associate category-specific lexical patterns with PCL in general, rather than learning a consistent notion of condescension across all categories.
\end{itemize}

% ============================================================
\clearpage
\section{Exercise 2: Exploratory Data Analysis (EDA)}
\label{sec:eda}

\subsection{EDA Technique 1: Class Distribution and Token Length Analysis}

\subsubsection*{Visual/Tabular Evidence}

\begin{figure}[H]
\centering
\includegraphics[width=\textwidth]{images/class_distribution_train_vs_dev.png}
\caption{Class distribution across train and dev splits. Both maintain the same $\sim$9.5:1 No-PCL:PCL ratio (train: 8,375 / 794 PCL; dev: 2,094 / 199 PCL).}
\label{fig:class_dist}
\end{figure}

\begin{figure}[H]
\centering
\includegraphics[width=\textwidth]{images/token_lengths_train_vs_dev.png}
\caption{Word-count distributions by class for train (left) and dev (right). PCL examples are slightly longer in both splits, with 95th percentiles below 200 words.}
\label{fig:token_len}
\end{figure}

\subsubsection*{Analysis}

The DPM corpus reflects a roughly \textbf{9.5:1 class imbalance} preserved uniformly across splits (train: 8,375 examples, 794 PCL; dev: 2,094 examples, 199 PCL), confirming the class ratio is consistent across both sets.

Token length distributions are also consistent across splits: PCL examples have a slightly higher mean word count ($\sim$54 words vs.\ $\sim$48 for No-PCL) with similar medians. Both splits share a right-skewed distribution; the 95th percentile is approximately 100--115 words across both classes and splits.

\subsubsection*{Impact Statement}

These observations directly inform two decisions applied in our pipeline:

\begin{enumerate}[noitemsep]
    \item \textbf{Weighted Random Sampler (WRS):} To counter the 9.5:1 imbalance, we sample training examples with per-class inverse-frequency weights, ensuring each mini-batch contains a roughly equal number of PCL and No-PCL examples. Without this, the gradient signal from the minority class would be overwhelmed and the model would converge to always predicting No-PCL.
    \item \textbf{\texttt{max\_length=256}:} We set the tokeniser's maximum sequence length to 256. Since the 95th percentile of sequence lengths is $\sim$115 words ($\sim$150 subword tokens), this ensures virtually no content is truncated while keeping memory usage feasible.
\end{enumerate}

\subsection{EDA Technique 2: Lexical Analysis (Log-Odds Contrast and Distinctive N-grams)}

\subsubsection*{Visual/Tabular Evidence}

\begin{figure}[H]
\centering
\includegraphics[width=0.78\textwidth]{images/log_odds_contrast.png}
\caption{Log-odds ratio of word frequencies in PCL vs.\ No-PCL examples (training set). Positive values indicate words more characteristic of PCL; negative values indicate words more characteristic of No-PCL.}
\label{fig:log_odds}
\end{figure}

\begin{figure}[H]
\centering
\includegraphics[width=\textwidth]{images/distinctive_bigrams.png}
\caption{Most distinctive bigrams in PCL vs.\ No-PCL examples (training set), measured by log-odds ratio.}
\label{fig:bigrams}
\end{figure}

\subsubsection*{Analysis}

The log-odds contrast plot shows that PCL examples disproportionately feature words relating to vulnerability and emotional framing: ``help'', ``need'', ``community'', ``support'', and ``people'' appear more often in PCL text. The bigram analysis identifies phrases such as ``in need'', ``help them'', ``these people'', and ``the homeless'' as distinctive of PCL, consistent with language that frames vulnerable communities as passive recipients of charity rather than active agents.

However, many of these words (e.g., ``help'', ``community'') are equally common in genuinely supportive, non-condescending writing. This is the central challenge of PCL detection: \emph{surface-level lexical cues alone are insufficient}. Two texts using almost identical vocabulary can differ fundamentally in tone. The patronising quality often lies in implied power dynamics, presuppositions about the group's capabilities, or the framing of their situation as hopeless, none of which are directly recoverable from word frequencies alone.

\subsubsection*{Impact Statement}

These findings motivate two decisions we apply throughout the pipeline:

\begin{enumerate}[noitemsep]
    \item \textbf{Pre-trained transformer models:} Since surface n-gram features are unreliable discriminators of PCL, we use RoBERTa-base [3] and DistilBERT [4] as our classifiers. Both are pre-trained on large corpora and encode contextual cues such as tone, framing, and presupposition that bag-of-words models cannot capture.

    \item \textbf{Metadata prepending:} We prepend the community keyword and country of origin to every input as structured context:\ \texttt{<e>keyword</e>\ <e>country</e>\ \{text\}}. This gives the model explicit prior knowledge about which vulnerable group is being discussed and the geopolitical setting, both of which influence the likelihood and style of PCL. The per-keyword analysis (Section~\ref{sec:local}) confirms that different keywords carry substantially different PCL base rates, validating this context signal.
\end{enumerate}

% ============================================================
\clearpage
\section{Exercise 3: Proposed Approach}
\label{sec:approach}

\subsection{Proposed Approach}

Our proposed system is a \textbf{dual-backbone ensemble} of \texttt{roberta-base} and \texttt{distilbert-base-\allowbreak{}uncased}, designed to address the key challenges from EDA: a 9.5:1 class imbalance, subtle pragmatic signals, and limited labelled positives.

\subsubsection*{Data Splits and Evaluation Protocol}

The official splits are used as follows:

\begin{itemize}[noitemsep]
    \item \textbf{Official train (8,375 examples):} split into an internal train set (90\%) and an internal validation set (10\%), both stratified by label. The validation set is used for early stopping, checkpoint selection, and threshold/weight tuning.
    \item \textbf{Official dev (2,094 examples):} treated as a held-out test set throughout development, since the official test set carries no labels. This allows us to measure generalisation on truly unseen data without leaking information into training decisions.
    \item \textbf{Official test (3,832 examples):} unlabelled; used only to generate the leaderboard submission files \texttt{dev.txt} and \texttt{test.txt}.
\end{itemize}

\subsubsection*{Input Representation}

Each input to both models is formatted as:
\begin{center}
    \texttt{<e>keyword</e>\ \ <e>country</e>\ \ \{paragraph text\}}
\end{center}
\begin{itemize}[noitemsep]
    \item HTML entities in the raw text are unescaped before tokenisation (e.g.\ \texttt{\&amp;} $\to$ \texttt{\&}).
    \item The community keyword and country are prepended as structured metadata to provide community-specific context at inference time.
    \item Inputs are tokenised with \texttt{max\_length=256} and truncation enabled (RoBERTa BPE tokeniser; DistilBERT WordPiece tokeniser).
\end{itemize}

\subsubsection*{Training Procedure}

Each backbone is trained independently for up to 8 epochs. An independent hyperparameter search is conducted for each architecture to identify the best-performing configuration:

\begin{itemize}[noitemsep]
    \item \textbf{Weighted Random Sampler (WRS):} per-class weights use a square-root-scaled inverse-frequency formula rather than full inverse frequency, providing a softer correction. Ambiguous positives (annotator agreement score 2 out of 4) receive an additional 0.5$\times$ down-weight to reduce noise from borderline labels.
    \item \textbf{Grouped layer-wise learning rate decay (LLRD):} transformer layers are divided into three groups; lower layers receive a smaller learning rate (base~$/ \lambda$) and upper layers a larger one (base~$\times \lambda$, with $\lambda = 1.6$). This preserves general pre-trained representations while allowing task-specific adaptation near the classifier head.
    \item \textbf{Cosine annealing with linear warmup:} the rate warms up linearly over the first 10\% of total training steps, then follows a cosine decay. This prevents large early updates from destabilising pre-trained weights.
    \item \textbf{Mid-training evaluation and checkpoint saving:} validation F1 is computed every 50 gradient update steps. Whenever a new best is reached, the model state is captured in memory. This fine-grained monitoring allows saving the best point within an epoch, not just at epoch boundaries.
    \item \textbf{Early stopping:} training halts after 15 consecutive evaluation checks (750 gradient steps) without improvement in validation F1.
    \item \textbf{Multi-seed selection:} each configuration is trained with three random seeds (42, 7, and 123). After all seeds complete, the checkpoint with the highest internal validation F1 is saved to disk. This reduces dependence on a single random initialisation.
\end{itemize}

\subsubsection*{Ensemble Weight Search and Prediction}

Once the best checkpoint for each backbone is identified, we combine them via weighted probability averaging:

\begin{itemize}[noitemsep]
    \item \textbf{2-D grid search:} a grid search is performed over the RoBERTa weight $w \in [0, 1]$ (in steps of 0.05) and the classification threshold $t \in [0.05, 0.95]$ (in steps of 0.01). The combined score for each example is $w \cdot p_{\text{RoBERTa}} + (1-w) \cdot p_{\text{DistilBERT}}$. The pair $(w^*, t^*)$ that maximises F1 on the internal validation set is selected. Since both backbones are trained with WRS, their probability outputs are well-calibrated and we expect the optimal weight to be close to 0.50, reflecting genuine complementarity rather than dominance by one model.
    \item \textbf{Automatic fallback:} if the best ensemble validation F1 does not exceed the solo RoBERTa validation F1, the ensemble is discarded and the solo model is used for submission. This prevents a miscalibrated DistilBERT from degrading performance.
    \item \textbf{Final predictions:} the selected system (ensemble or solo) is applied to both the official dev and official test sets using $(w^*, t^*)$, producing \texttt{dev.txt} and \texttt{test.txt}.
\end{itemize}

\subsection{Rationale and Expected Outcome}

\subsubsection*{Base Models: RoBERTa-base and DistilBERT}

Standard BERT trained on 16~GB of text using masked-language modelling (MLM) and next-sentence prediction (NSP). RoBERTa [3] improved on BERT by training for longer on 160~GB of text, removing NSP, and switching to dynamic masking, producing representations that are more robust to domain shift. DistilBERT [4] compresses BERT via knowledge distillation into a model with 40\% fewer parameters and 60\% faster inference, retaining $\sim$97\% of BERT-level performance.

We choose these two models specifically because their differences make them complementary: distinct architectures (12 vs.\ 6 layers), tokenisers (byte-pair encoding vs.\ WordPiece), and pre-training signals result in partially uncorrelated errors. This is the key condition for an ensemble gain: when one model misclassifies an example the other does not, averaging their outputs corrects the error. For a task like PCL, where the signal is subtle and the dataset small, combining two well-chosen backbones is expected to be more effective than scaling a single model.

\subsubsection*{Training Methodology}

Each of the following components addresses a specific challenge identified in the EDA or known from the fine-tuning literature:

\textbf{Weighted Random Sampler (WRS):} Addresses two problems simultaneously. First, it corrects the 9.5:1 class imbalance by ensuring balanced mini-batches, preventing convergence to the degenerate always-predict-No-PCL solution. Second, and more critically, it is essential for probability calibration: without WRS, the model sees $\sim$10$\times$ more negative examples per batch and learns to push all outputs toward zero, forcing the optimal threshold as low as $t=0.07$--$0.14$ in our experiments, a value specific to the validation distribution that does not transfer to unseen data.

\textbf{Grouped Layer-wise Learning Rate Decay (LLRD):} Prevents catastrophic forgetting. Updating all layers at the same rate risks overwriting the general linguistic representations learnt during pre-training. We assign smaller learning rates to lower transformer layers ($2 \times 10^{-5}$) and progressively larger rates to upper layers and the classifier head, preserving general knowledge while allowing task-specific adaptation near the output.

\textbf{Cosine annealing with linear warmup:} The warmup phase prevents large early-stage updates from destabilising pre-trained weights. The subsequent cosine decay reduces the rate smoothly, favouring convergence to a flatter and more generalisable minimum than step-decay or constant schedules.

\textbf{Post-hoc threshold tuning:} Even with WRS, the F1-optimal boundary for individual models is rarely 0.50 (the best solo RoBERTa requires $t=0.75$; seeds span 0.12--0.87). We sweep the threshold on the held-out validation set to find the F1-maximising value without using test data.

\textbf{Multi-seed training:} Fine-tuning large transformers on small, imbalanced datasets is unstable. The same configuration with three seeds can produce F1 values spanning $\sim$0.03 and optimal thresholds spanning 0.75. Training three seeds and selecting the best checkpoint avoids over-reliance on a single random initialisation.

\textbf{Metadata prepending:} The log-odds analysis (Section~\ref{sec:eda}) shows that PCL frequency varies substantially across keywords. By prepending the community keyword and country to every input, the model receives explicit context about which vulnerable group is being discussed, helping it calibrate appropriately per category.

\subsubsection*{Overall Expected Outcomes and Potential Drawbacks}

Taken together, these components prioritise test-set generalisation over training performance. WRS and LLRD act as implicit regularisers that reduce overfitting to the training distribution; multi-seed selection avoids committing to a lucky or unlucky random initialisation; and ensembling two architecturally distinct, WRS-calibrated backbones smooths seed-specific biases and combines complementary error patterns. We expect the ensemble to reach F1~$\sim$0.55--0.62 on the dev set, with probability outputs compatible at a shared threshold of 0.50 and equal weights, without requiring per-model calibration.

There are genuine trade-offs. WRS over-samples positive examples, which may amplify noise from borderline or ambiguously labelled PCL instances and slow convergence. LLRD group boundaries are heuristic; they may not align precisely with how the model internally partitions general vs.\ task-specific knowledge. Tripling the compute for multi-seed training reduces the number of hyperparameter configurations we can explore. Finally, ensembling makes the system harder to interpret and roughly doubles inference time. The ablation study in Section~\ref{sec:ablation} tests these expectations and trade-offs systematically.

% ============================================================
\clearpage
\section{Exercise 4 \& 5.1: Repository and Global Evaluation}

\textbf{BestModel directory.} The \texttt{BestModel/} folder at the repository root contains all artefacts required by Exercise~4:

\begin{itemize}[noitemsep]
    \item \texttt{model.ipynb}: the full training pipeline (data loading, WRS, LLRD, multi-seed loop, threshold tuning, ensemble, evaluation) as proposed in Exercise~3.
    \item \texttt{roberta-base/}: saved weights for the best RoBERTa-base checkpoint (drop=0.0, bs=4, ga=4, seed=7).
    \item \texttt{distilbert/}: saved weights for the best DistilBERT checkpoint (drop=0.3, bs=8, ga=2, seed=7).
\end{itemize}

\textbf{Submission files (Exercise 5.1).} The files \texttt{dev.txt} and \texttt{test.txt} are in the repository root, each containing one binary prediction per line (0 = No PCL, 1 = PCL):

\begin{itemize}[noitemsep]
    \item \texttt{dev.txt}: 2,094 lines, matching the official dev set exactly.
    \item \texttt{test.txt}: 3,832 lines, matching the official test set exactly.
\end{itemize}

The ensemble achieves \textbf{F1 = 0.6069} on the official dev set, surpassing the 0.48 baseline by $+$0.1269. The GitHub repository link is on the front page of this report.

% ============================================================
\clearpage
\section{Exercise 5: Local Evaluation}
\label{sec:local}

This section presents the local evaluation of our submitted system. Table~\ref{tab:final_results} summarises results across all trained configurations; the subsections below analyse performance in detail.

\begin{table}[H]
\centering
\begin{tabular}{lccc}
\toprule
System & Val F1 & Dev F1 & Threshold \\
\midrule
Solo RoBERTa (drop=0.0, bs=4, ga=4, seed=7)       & 0.5795 & 0.5955 & 0.75 \\
Solo DistilBERT (drop=0.3, bs=8, ga=2, seed=7)     & 0.5818 & 0.5463 & 0.50 \\
\textbf{Ensemble (w=0.50/0.50)}                    & \textbf{0.5882} & \textbf{0.6069} & \textbf{0.50} \\
Official baseline (RoBERTa-base, no tuning)        & -- & 0.48 & -- \\
\bottomrule
\end{tabular}
\caption{Final system comparison on the official dev set. Val F1 is the best-threshold score on the held-out validation split.}
\label{tab:final_results}
\end{table}

% ============================================================
\subsection{Quantitative Performance Evaluation}

\subsubsection*{Confusion Matrix}

\begin{figure}[H]
\centering
\includegraphics[width=0.88\textwidth]{images/confusion_matrix.png}
\caption{Confusion matrix on the dev set (counts, left; row-normalised recall per class, right). Ensemble weights 0.50/0.50, threshold 0.50.}
\label{fig:cm}
\end{figure}

\begin{table}[H]
\centering
\begin{tabular}{lcc}
\toprule
Metric & Value & Interpretation \\
\midrule
True Negatives (TN)    & 1,793 & Correct No-PCL predictions \\
False Positives (FP)   & 102   & No-PCL flagged as PCL \\
False Negatives (FN)   & 69    & PCL examples missed \\
True Positives (TP)    & 130   & Correct PCL predictions \\
\midrule
Precision (PCL)        & 0.560 & 56\% of flagged examples are truly PCL \\
Recall (PCL)           & 0.653 & 65.3\% of all PCL examples detected \\
F1 (PCL)               & 0.603 & Harmonic mean of precision and recall \\
False Positive Rate    & 0.054 & 5.4\% of No-PCL examples incorrectly flagged \\
False Negative Rate    & 0.347 & 34.7\% of PCL examples missed \\
\bottomrule
\end{tabular}
\caption{Confusion matrix statistics for the ensemble on the dev set.}
\label{tab:cm_stats}
\end{table}

The ensemble operates at a precision-oriented point: it correctly identifies 130 of 199 PCL examples (recall 0.653) but misses 69, and incorrectly flags only 102 No-PCL examples (FPR 5.4\%). The asymmetry between FPR (5.4\%) and FNR (34.7\%) reflects a cautious decision boundary: the model predicts PCL only when the averaged ensemble probability clears 0.50, which, thanks to WRS calibration, functions as a genuine midpoint rather than an artefact of probability collapse. Most predictions default to No-PCL; PCL is flagged only when both backbones agree.

This behaviour is expected given the 9.5:1 imbalance: even with balanced mini-batches during training, the test distribution remains skewed and drives the model toward the majority class on borderline examples. The F1 of 0.603 reflects a reasonable trade-off for the SemEval criterion: lowering the threshold would recover more PCL cases but at the cost of more FPs, which are themselves misleading in a content-flagging application.

\subsubsection*{Precision-Recall Curve}

\begin{figure}[H]
\centering
\includegraphics[width=0.72\textwidth]{images/pr_curve_threshold.png}
\caption{Precision--Recall curve for the ensemble on the dev set. AUC-PR = 0.5906. The dashed vertical line marks the chosen threshold (0.50).}
\label{fig:pr_curve}
\end{figure}

The AUC-PR of \textbf{0.5906} is roughly 6$\times$ higher than the no-skill baseline of $\approx$0.095 (equal to the PCL base rate), confirming that the ensemble has learnt genuine discriminative signal rather than class-frequency heuristics.

The shape of the curve reveals three distinct operating regimes, each appropriate for a different use case:

\begin{itemize}[noitemsep]
    \item \textbf{High recall (threshold $\lesssim$0.25):} recall $>$0.85 at precision $\sim$0.25--0.30. Suitable for auditing or research where missing PCL is the primary concern and false positives can be reviewed manually.
    \item \textbf{Balanced (threshold $\approx$0.50, chosen):} precision = 0.560, recall = 0.653, F1 = 0.603. Selected in line with the SemEval F1 criterion.
    \item \textbf{High precision (threshold $\approx$0.75):} precision $>$0.70 at recall $\approx$0.50. Preferred when false positives are costly, for example in automated content moderation where incorrect flags cause direct harm.
\end{itemize}

The steep drop in precision at high recall (left side of the curve) is characteristic of imbalanced classification: only a small subset of examples carry strong PCL signal, and model confidence degrades quickly as more borderline cases are included. This is consistent with the per-keyword findings in Section~\ref{sec:errors}: categories with distinctive lexical markers (\textit{in-need}, \textit{homeless}) contribute most of the high-confidence TP detections, while subtler categories reduce precision as the threshold is lowered.

\subsubsection*{Training Stability: Multi-Seed Analysis}

\begin{figure}[H]
\centering
\includegraphics[width=0.70\textwidth]{images/multi_seed_comparison.png}
\caption{RoBERTa-base (drop=0.0, bs=4, ga=4) trained with three different random seeds. Bar heights show best-threshold F1 on the dev set.}
\label{fig:multi_seed}
\end{figure}

\begin{table}[H]
\centering
\begin{tabular}{lccc}
\toprule
Seed & Best threshold & Dev F1 & AUC-PR \\
\midrule
7   & 0.87 & 0.6034 & 0.5835 \\
42  & 0.12 & 0.5720 & 0.5792 \\
123 & 0.62 & 0.5693 & 0.5666 \\
\bottomrule
\end{tabular}
\caption{RoBERTa-base multi-seed results on the dev set (per-seed threshold tuning applied).}
\label{tab:multi_seed}
\end{table}

Training the same configuration with three random seeds reveals two important observations.

First, the dev F1 scores span 0.569--0.603 (a gap of 0.034). Given that our target is to improve over the 0.48 baseline, this variance is non-trivial: there is a 2-in-3 chance of selecting a weaker checkpoint without multi-seed evaluation. Post-hoc threshold tuning is essential to extract each seed's best performance; defaulting to $t=0.50$ would be sub-optimal for two of the three seeds.

Second, and more telling, are the optimal thresholds: 0.12, 0.62, and 0.87 for seeds 42, 123, and 7 respectively, a spread of 0.75. Despite this, the AUC-PR values are tightly clustered at 0.567--0.584 (gap of only 0.017). This contrast is informative: similar AUC-PR means the three checkpoints have learnt a similar underlying discriminative signal, but their output probabilities are on very different scales. Seed 42 collapses probabilities toward zero (requiring $t=0.12$ to recover F1, a sign of residual miscalibration even with WRS); seed 7 distributes probabilities more uniformly across $[0,1]$ and demands a high threshold to avoid over-flagging. This is the core motivation for multi-seed training: not just variance in F1, but variance in calibration that would make naive 50/50 ensembling unreliable without WRS.

The ensemble (dev F1 = 0.6069) outperforms the best single seed (0.6034, seed~7) by combining independently calibrated probability outputs. With both backbones trained with WRS, the averaged probabilities are compatible at a single threshold of 0.50, eliminating the need for per-model calibration.

% ============================================================
\clearpage
\subsection{Ablation Study: Verifying Design Decisions}
\label{sec:ablation}

We conducted a three-part ablation to verify the design decisions proposed in Section~3. All controlled experiments use \texttt{distilbert-base-uncased} (seed=42) as the backbone, a faster and smaller model that enables rapid iteration. Results confirm that the submitted ensemble achieves the highest test F1 (0.6069) and is the only configuration that generalises from validation to test.

\subsubsection*{Part A: Component Ablation}

Each of the five training-time components (metadata prepending, WRS, grouped LLRD, linear warmup, threshold tuning) was removed one at a time while keeping all other settings fixed. All variants use up to 10 epochs with early stopping (patience=3). Full code is in \texttt{5\_evaluation/ablation\_study.ipynb}.

\begin{table}[H]
\centering
\begin{tabular}{lcccc}
\toprule
Variant & Val F1 & Val threshold & Test F1 & $\Delta$ Test F1 \\
\midrule
Full model (all components ON) & 0.5426 & 0.67 & 0.5031 & 0.000 \\
w/o Metadata prepending        & 0.5699 & 0.62 & 0.5367 & $+$0.034 \\
w/o Weighted Random Sampler    & 0.5882 & 0.14 & 0.5458 & $+$0.043 \\
w/o Grouped LLRD               & 0.5333 & 0.78 & 0.5284 & $+$0.025 \\
w/o Linear warmup              & 0.5611 & 0.56 & 0.5431 & $+$0.040 \\
Fixed threshold $t=0.50$       & 0.5022 & 0.50 & 0.4893 & $-$0.014 \\
\midrule
RoBERTa drop=0.0 (best, from log) & 0.5979 & 0.87 & 0.6034 & $+$0.100 \\
\textbf{Ensemble (submitted)}     & \textbf{0.5882} & \textbf{0.50} & \textbf{0.6069} & $+$0.104 \\
\bottomrule
\end{tabular}
\caption{Component-level ablation (DistilBERT, seed=42). $\Delta$ Test F1 = test F1 minus the DistilBERT full-model baseline (0.5031); positive means performance \emph{improved} when the component was removed, which is counter-intuitive and explained in the text.}
\label{tab:ablation_component}
\end{table}

\begin{figure}[H]
\centering
\includegraphics[width=0.75\textwidth]{images/ablation_bar.png}
\caption{Test F1 for each ablation variant and the final submitted models (DistilBERT full model baseline: 0.5031).}
\label{fig:ablation_bar}
\end{figure}

The positive deltas (removing a component \emph{raises} test F1) are counter-intuitive but confirm the calibration argument made in Section~3. As a standalone system, DistilBERT with all components active is already miscalibrated: it requires $t=0.67$, well above the neutral 0.50. Removing WRS worsens this to $t=0.14$, confirming that WRS removal pushes probabilities toward zero as predicted. The low threshold recovers validation F1 but does not transfer to the test set, consistent with our generalisation concern.

The only component with a negative delta is threshold tuning: fixing $t=0.50$ drops F1 by $-$0.014, confirming that val-set threshold sweeping is necessary for individual models. In the full ensemble, both backbones are trained with WRS, their calibrated probabilities average naturally at $t=0.50$, and test F1 = 0.6069, over 0.10 above the standalone DistilBERT baseline, exactly as anticipated.

\subsubsection*{Part B: Hyperparameter Search (Dropout and Batch Size)}

\begin{table}[H]
\centering
\begin{tabular}{llccc}
\toprule
Backbone & Config & seed=42 & seed=7 & seed=123 \\
\midrule
RoBERTa-base & drop=0.0, bs=8, ga=2  & 0.5778 & 0.5786 & \textbf{0.5979} \\
RoBERTa-base & drop=0.0, bs=4, ga=4  & 0.5890 & 0.5795 & \textbf{0.5979} \\
RoBERTa-base & drop=0.5, bs=8, ga=2  & 0.2471 & \textbf{0.2768} & 0.1753 \\
RoBERTa-base & drop=0.5, bs=4, ga=4  & 0.3000 & 0.2249 & \textbf{0.3203} \\
\midrule
DistilBERT   & drop=0.1, bs=8, ga=2  & 0.5596 & \textbf{0.5965} & 0.5965 \\
DistilBERT   & drop=0.1, bs=4, ga=4  & 0.5761 & 0.5829 & \textbf{0.5889} \\
DistilBERT   & drop=0.3, bs=8, ga=2  & 0.5775 & \textbf{0.5818} & 0.5657 \\
\midrule
\textbf{Ensemble} & best RoBERTa + best DistilBERT & \multicolumn{3}{c}{\textbf{Dev F1 = 0.6069}} \\
\bottomrule
\end{tabular}
\caption{Hyperparameter search: best-threshold validation F1 per seed. \textit{bs}=batch size, \textit{ga}=gradient accumulation steps (effective batch = bs\,$\times$\,ga).}
\label{tab:ablation_dropout}
\end{table}

Three findings confirm the Section~3 expectations. \textbf{(1) Dropout is the most destructive hyperparameter:} RoBERTa with drop=0.5 collapses from val F1~$\sim$0.598 to $\sim$0.28--0.32, consistent with the gradient-starvation prediction; drop=0.0 is optimal for RoBERTa while DistilBERT tolerates modest dropout (0.1--0.3) due to its smaller capacity. \textbf{(2) Effective batch size is what matters:} both configurations (bs=8/ga=2 and bs=4/ga=4, both effective batch=16) reach the same peak val F1, confirming that gradient quality rather than update frequency drives performance. \textbf{(3) Ensembling provides the largest single gain:} the best solo RoBERTa (dev F1=0.6034) is outperformed by the ensemble (dev F1=0.6069), with Figure~\ref{fig:rvd} directly showing the partially uncorrelated error sets that enable this.

\begin{figure}[H]
\centering
\includegraphics[width=0.80\textwidth]{images/roberta_vs_distilbert.png}
\caption{Error overlap between the best RoBERTa and DistilBERT checkpoints on the dev set. Predictions unique to one model are the primary motivation for ensembling.}
\label{fig:rvd}
\end{figure}

\subsubsection*{Part C: Generalisation Verification}

Based on the component ablation results, we trained two new DistilBERT variants across three seeds (selecting the best checkpoint per variant), specifically designed to capitalise on the positive ablation deltas observed above:

\begin{itemize}[noitemsep]
    \item \textbf{Variant A}: all four components (metadata, WRS, LLRD, warmup) removed; dropout=0.3.
    \item \textbf{Variant B}: only WRS and warmup removed; dropout=0.1.
\end{itemize}

Both variants achieved higher validation F1 during training (Variant A: 0.5851; Variant B: 0.5868) compared to the submitted DistilBERT checkpoint, which motivated evaluating them as potential ensemble partners.

\begin{table}[H]
\centering
\begin{tabular}{lcccc}
\toprule
Model & Val F1* & Test F1 & Val$\to$Test gap & $t$ \\
\midrule
RoBERTa (submitted)    & --     & 0.5955 & --       & 0.75 \\
DistilBERT (submitted) & --     & 0.5455 & --       & 0.53 \\
Variant A (seed=123)   & 0.5851 & 0.5628 & $-$0.022 & 0.07 \\
Variant B (seed=7)     & 0.5868 & 0.5208 & $-$0.066 & 0.09 \\
\midrule
\textbf{Submitted ensemble}   & \textbf{0.5882} & \textbf{0.6069} & \textbf{$+$0.019} & \textbf{0.50} \\
Ensemble + Variant A          & 0.6115 & 0.5854 & $-$0.026 & 0.29 \\
Ensemble + Variant B          & 0.6207 & 0.5753 & $-$0.045 & 0.26 \\
\bottomrule
\end{tabular}
\caption{Generalisation verification. Val F1* is the best-threshold score on the internal validation split. The submitted ensemble is the only system where test F1 exceeds val F1 ($+$0.019), confirming genuine generalisation.}
\label{tab:gen_verify}
\end{table}

Despite higher training-time validation scores, both new variants generalise poorly: val F1 up to 0.6207 but test F1 as low as 0.5753. The root cause is probability miscalibration from removing WRS. Without WRS, the model sees $\sim$10$\times$ more negative examples per batch and pushes all outputs toward zero, requiring extreme decision thresholds ($t=0.07$ and $t=0.09$) and heavily skewed ensemble weights (25\%/75\%) that are specific to the validation set and do not transfer to the test set. The submitted ensemble, trained with WRS in both backbones, is well-calibrated at $t=0.50$ with equal weights and is the only system that improves from validation to test.

\textbf{Key conclusion:} WRS is not just a class-imbalance corrector; it is essential for probability calibration and generalisation. Higher validation F1 achieved by removing it is an artefact of aggressive threshold tuning, not genuine improvement.

% ============================================================
\clearpage
\subsection{Error Analysis: Failure Modes and Per-Keyword Performance}
\label{sec:errors}

We compare the ensemble against the solo RoBERTa model as a reference system, since both sets of predictions are available.

\subsubsection*{Comparison with Solo RoBERTa (Baseline Reference)}

The solo RoBERTa model achieves dev F1 = 0.5955; the ensemble achieves 0.6069, a gain of $+$0.011. This improvement comes from cases where RoBERTa errs but DistilBERT does not, without introducing proportionally many new errors. Figure~\ref{fig:rvd} (Section~\ref{sec:ablation}) visualises the four prediction overlap categories:

\begin{itemize}[noitemsep]
    \item \textbf{(a) Both correct:} the large majority of examples; both models agree on the dominant No-PCL class.
    \item \textbf{(b) Both wrong:} hard PCL cases where neither model has enough signal. These are typically indirect, politically framed, or presupposition-based PCL with no surface-level markers.
    \item \textbf{(c) RoBERTa correct, ensemble wrong:} rare; the ensemble can inherit DistilBERT's errors if DistilBERT is confidently wrong and RoBERTa is uncertain. This happens more on low-frequency categories such as \textit{migrant} and \textit{women}.
    \item \textbf{(d) Ensemble correct, RoBERTa wrong:} the primary source of the ensemble's gain. DistilBERT's different inductive biases (lighter architecture, distinct tokenisation) allow it to catch some PCL examples that RoBERTa misses, particularly in categories such as \textit{in-need} and \textit{homeless} where PCL is more lexically distinctive.
\end{itemize}

The asymmetry (more examples in category (d) than (c)) confirms that the ensemble genuinely diversifies predictions rather than averaging noise.

\subsubsection*{Manual Inspection of False Positives and False Negatives}

\begin{table}[H]
\centering
\small
\begin{tabular}{p{0.8cm}p{1.5cm}p{1.2cm}p{6.8cm}}
\toprule
Type & Keyword & Country & Text (truncated) \\
\midrule
FP & in-need & ZA & \textit{``Adopt a Mission serves as a platform for the church and for likeminded people to reach out to unemployed families in the communities\ldots''} \\[4pt]
FP & poor-families & KE & \textit{``We have identified extremely poor families, orphaned children, the elderly and those living with HIV/AIDS but with no source of\ldots''} \\
\midrule
FN & disabled & HK & \textit{``Cheung said 20 disabled undergraduate students from seven universities will start their eight-week internship in government departments\ldots''} \\[4pt]
FN & poor-families & AU & \textit{``Michael Gove's recent suggestion that inadequate financial management skills among poor families are to blame for the increasing\ldots''} \\[4pt]
FN & hopeless & BD & \textit{``Calling for an immediate political solution to resolve the crisis in Syria, Erdogan said: `Turkey's incursion into northern Syria\ldots'\thinspace''} \\
\bottomrule
\end{tabular}
\caption{Representative misclassifications: false positives (top) and false negatives (bottom).}
\label{tab:errors}
\end{table}

\textbf{False Positives.} Both FP examples use factual or organisational language about vulnerable communities. The model latches onto vulnerability-associated vocabulary (``poor'', ``unemployed'', ``HIV/AIDS'') rather than the framing or power dynamics that constitute genuine PCL.

\textbf{False Negatives.} The FN examples involve implicit PCL. The ``Michael Gove'' example is a \emph{Presupposition} PCL case: it implies poor families are financially incompetent through bureaucratic, indirect language. The ``Syria/hopeless'' example embeds PCL in political framing. Both require pragmatic reasoning and world knowledge that go beyond surface token patterns, representing the fundamental limit of current transformer classifiers on this task.

\subsubsection*{Per-Keyword Analysis}

\begin{figure}[H]
\centering
\includegraphics[width=0.85\textwidth]{images/per_keyword_analysis.png}
\caption{F1, precision, and recall on the dev set broken down by vulnerable-community keyword.}
\label{fig:per_keyword}
\end{figure}

\begin{table}[H]
\centering
\begin{tabular}{lcccc}
\toprule
Keyword & PCL Rate & Precision & Recall & F1 \\
\midrule
in-need       & 14.6\% & 0.596 & 0.939 & 0.729 \\
homeless      & 13.7\% & 0.656 & 0.724 & 0.689 \\
immigrant     & 3.2\%  & 1.000 & 0.429 & 0.600 \\
poor-families & 20.0\% & 0.579 & 0.579 & 0.579 \\
vulnerable    & 9.6\%  & 0.520 & 0.650 & 0.578 \\
disabled      & 7.2\%  & 0.500 & 0.571 & 0.533 \\
hopeless      & 12.0\% & 0.447 & 0.654 & 0.531 \\
refugee       & 6.9\%  & 0.500 & 0.538 & 0.519 \\
migrant       & 2.4\%  & 0.667 & 0.400 & 0.500 \\
women         & 6.0\%  & 0.545 & 0.429 & 0.480 \\
\bottomrule
\end{tabular}
\caption{Per-keyword performance on the dev set. PCL Rate = fraction of examples labelled PCL for that keyword.}
\label{tab:per_keyword}
\end{table}

Performance correlates with the PCL base rate per keyword: \textbf{in-need} (F1=0.729, rate=14.6\%) and \textbf{homeless} (F1=0.689, rate=13.7\%) are the strongest-performing categories, while \textbf{women} (F1=0.480, rate=6.0\%), \textbf{migrant} (F1=0.500, rate=2.4\%), and \textbf{refugee} (F1=0.519, rate=6.9\%) are the weakest. This pattern suggests the model exploits keyword-specific lexical cues rather than a fully transferable notion of PCL. \textbf{Immigrant} is a notable outlier: perfect precision (1.000) but low recall (0.429). Every flagged case is correct, but most PCL examples in this category are missed, likely because immigrant-related PCL relies on subtle implication rather than direct lexical markers.

% ============================================================
\clearpage
\section{Discussion and Conclusion}
\label{sec:discussion}

This report has described a pipeline for binary PCL detection on SemEval 2022 Task 4, Subtask 1. EDA revealed a 9.5:1 class imbalance and that PCL is characterised by pragmatic signals rather than surface keywords alone, motivating the use of pre-trained transformers and a Weighted Random Sampler. We proposed and implemented a dual-backbone ensemble (RoBERTa-base + DistilBERT) trained with grouped LLRD, cosine annealing, metadata prepending, multi-seed training, and threshold tuning, achieving dev set F1 = 0.6069 ($+$0.1269 above the 0.48 official baseline).

The ablation study (Section~\ref{sec:ablation}) confirms that WRS is the most critical component for generalisation: removing it raises validation F1 through aggressive threshold tuning, but this does not transfer to unseen data. Classifier-head dropout is the most destructive hyperparameter in the grid search, and ensembling two architecturally distinct backbones provides a consistent gain over any individual checkpoint.

The model still struggles with indirect, politically framed, or presupposition-heavy PCL cases that require pragmatic reasoning beyond surface token patterns (Section~\ref{sec:errors}). Per-keyword performance correlates with the PCL base rate, suggesting the model captures some keyword-specific patterns rather than a general notion of PCL.

With more time and compute, the directions most likely to yield gains are: (1) \textbf{larger models} such as \texttt{deberta-v3-base}, which have shown strong results on comparable tasks; (2) \textbf{data augmentation} via back-translation to increase positive training examples; (3) \textbf{multi-task learning} using the Subtask 2 category labels as auxiliary supervision to sharpen the binary classifier.

% ============================================================
\clearpage
\section*{References}

\begin{enumerate}[label={[\arabic*]}]
    \item Pérez-Almendros, C., Espinosa-Anke, L., \& Schockaert, S. (2020). \textit{Don't Patronize Me! An Annotated Dataset with Patronizing and Condescending Language towards Vulnerable Communities}. In \emph{Proceedings of COLING 2020}, pp.\ 5891--5902.

    \item Pérez-Almendros, C., Espinosa-Anke, L., \& Schockaert, S. (2022). \textit{SemEval-2022 Task 4: Patronizing and Condescending Language Detection}. In \emph{Proceedings of SemEval 2022}, pp.\ 1--11.

    \item Liu, Y., Ott, M., Goyal, N., Du, J., Joshi, M., Chen, D., Levy, O., Lewis, M., Zettlemoyer, L., \& Stoyanov, V. (2019). \textit{RoBERTa: A Robustly Optimized BERT Pretraining Approach}. arXiv:1907.11692.

    \item Sanh, V., Debut, L., Chaumond, J., \& Wolf, T. (2019). \textit{DistilBERT, a distilled version of BERT: smaller, faster, cheaper and lighter}. arXiv:1910.01108.
\end{enumerate}

\end{document}
